\documentclass[11pt]{article}
\usepackage[T1]{fontenc}
\usepackage{textcomp}
\usepackage[margin=0.75in]{geometry}
\usepackage{amsmath,amssymb,amsthm}
\usepackage{array,booktabs}
\usepackage{hyperref}
\setlength{\parindent}{0pt}
\newtheorem{theorem}{Theorem}
\theoremstyle{definition}
\newtheorem{definition}{Definition}
\title{Flow Proof Artifact: Comb-choose-sym}
\author{Generated by \texttt{flow doc proof}}
\date{Flow Proof Book}
\begin{document}
\maketitle
Binomial coefficients are symmetric in k and n-k.\\[0.5em]
\textbf{Source.} graham-knuth-patashnik — *Concrete Mathematics*\\[0.5em]
\bigskip
\noindent{\large \textbf{Derived fact 1.} \textit{choose(n, k) = choose(n, n - k)} \hfill Comb \cdot choose \cdot symmetry in k and n minus k}\par
\smallskip\noindent\textit{Coordinate: Comb \cdot choose \cdot symmetry in k and n minus k}\par
\smallskip\noindent\textit{Source: graham-knuth-patashnik}\par
\smallskip\noindent\textit{Built on: pascal recurrence holds, for choose on Comb, choosing none gives one, for choose on Comb}\par
\medskip
\noindent\textbf{Goal.} choose(n, k) = choose(n, n - k)\par
\noindent$\displaystyle \forall n \in \mathbb{N} \forall k \in \mathbb{N}\quad choose(n, k) = choose(n, n - k)$\par
\medskip
\noindent
\begin{tabular}{@{} >{\raggedright\arraybackslash}p{0.06\textwidth} >{\raggedright\arraybackslash}p{0.47\textwidth} >{\raggedleft\arraybackslash}p{0.41\textwidth} @{}}
\textbf{\#} & \textbf{Proof} & \textbf{Mathematics} \\
\midrule
\textbf{1.} & We prove that symmetry in k and n minus k for choose on Comb. & \\[0.45em]
\textbf{2.} & We proceed by induction on n: first the base case, then the inductive step. & \\[0.45em]
\textbf{3.} & Consider the base case in which k is zero. & \\[0.45em]
\textbf{4.} & We invoke the definitional clause governing choose on Comb: choosing none gives one, for choose on Comb (instantiated for n). & \\[0.45em]
\textbf{5.} & From step 3 and step 4, we can deduce that choose(n, k) equals choose(n, n - k). This establishes the base case (see step 3 and step 4). Hence proven. & $\displaystyle choose(n, k) = choose(n, n - k)$ \\[0.45em]
\textbf{6.} & For the inductive step, suppose the claim holds for all smaller values. & \\[0.45em]
\textbf{7.} & Under the supposition in step 6, let j denote the predecessor of k. & \\[0.45em]
\textbf{8.} & Under the supposition in step 6, we cross the inductive boundary: assume the claim holds for pred(n) (the induction hypothesis). & $\displaystyle choose(pred(n), j) = choose(pred(n), pred(n) - j)$ \\[0.45em]
\textbf{9.} & We invoke the definitional clause governing choose on Comb: pascal recurrence holds, for choose on Comb (instantiated for n, k). & \\[0.45em]
\textbf{10.} & From step 6, step 7, step 8, and step 9, this implies choose(n, k) equals choose(n, n - k). Hence proven. & $\displaystyle choose(n, k) = choose(n, n - k)$ \\[0.45em]
\bottomrule
\end{tabular}
\label{thm:Comb-choose-symmetry_in_k_and_n_minus_k}
\par\medskip\hrule
\end{document}
